\documentclass[11pt]{article}
\usepackage[margin=1in]{geometry}
\usepackage{booktabs}
\usepackage{longtable}
\usepackage{hyperref}
\usepackage{xcolor}
\usepackage{tabularx}

\title{Independent Replication of McCubbin et al.\ 2020:\\
A Pan-Genome Guided Metabolic Network Reconstruction of\\
Five \emph{Propionibacterium} Species}
\author{Independent Replication Project (BVBRC-19)}
\date{2026-06-27}

\begin{document}
\maketitle

\begin{abstract}
We independently replicate the three numerical pillars of McCubbin et al.\ 2020
(\emph{Genes} 11:1115, PMC7650540): (i) inter-species pan-genome reconstruction,
(ii) species-distinguishing pathway presence/absence, and (iii) genome-scale
flux-balance behavior of the six published genome-scale metabolic models (GEMs).
Working from the authors' deposited Supplementary File 4 GenBank files (17{,}525
CDS with translations across six inter-species representative strains) we recover
\textbf{909 core clusters} against the paper's stated \textbf{792--906} range
(within 0.3\% of the upper bound), reproduce the \textbf{open} pan-genome
accumulation shape ($+438$--$698$ new clusters per added genome), and confirm
\textbf{14/14} presence/absence pathway claims exactly, including the diagnostic
absences (transaldolase absent only in \emph{P.\ avidum}; sucrose-6-P hydrolase
only in \emph{P.\ acidipropionici} + \emph{P.\ propionicum}). All six deposited
GEMs solve to positive growth with propionate as the major fermentation product
and reproduce the reported auxotrophy hierarchy. Verdict: \textbf{PARTIAL} --
downstream analyses replicate cleanly, but the upstream KBase/RAST/GLIMMER
re-annotation step is \emph{not} independently re-derived (we start from the
authors' deposited standardized GenBank files, not raw NCBI assemblies).
\end{abstract}

\section{Paper Under Test}
McCubbin, T., \emph{et al.}\ (2020). \emph{A Pan-Genome Guided Metabolic Network
Reconstruction of Five Propionibacterium Species Reveals Extensive Metabolic
Diversity.} \emph{Genes} 11(10):1115. PMC7650540.

The paper reports 6 representative GEMs plus pan-genomic comparison across
6 inter-species representative strains drawn from 16 closed NCBI
\emph{Propionibacterium} genomes.

\section{Genomes Analyzed}
\begin{center}
\begin{tabular}{lll}
\toprule
Tag & Strain & CDS w/ translation (this rerun) \\
\midrule
PAC\_4875  & \emph{P.\ acidipropionici} ATCC 4875 & 3{,}365 \\
PAC\_55737 & \emph{P.\ acidipropionici} ATCC 55737 & 3{,}512 \\
PSHE       & \emph{P.\ freudenreichii} subsp.\ \emph{shermanii} CIRM-BIA1 & 2{,}406 \\
PAVI       & \emph{P.\ avidum} 44067 & 2{,}449 \\
PACN       & \emph{P.\ acnes} 6609 & 2{,}520 \\
PPRO       & \emph{P.\ propionicum} F0230a & 3{,}273 \\
\midrule
\textbf{Total} & & \textbf{17{,}525} \\
\bottomrule
\end{tabular}
\end{center}

\section{Layer 1 --- Pan-genome Reconstruction}
Pipeline: extract \texttt{/translation=} proteins from each GBK $\rightarrow$
all-vs-all \texttt{blastp} ($17{,}525 \times 17{,}525$, e-value $\le 10^{-5}$)
$\rightarrow$ filter to $\ge 30\%$ identity and $\ge 75\%$ coverage of the
shorter protein (paper's OrthoMCL 75\% coverage) $\rightarrow$ symmetrize edges
$\rightarrow$ bit-score-weighted graph $\rightarrow$ \textbf{MCL inflation 1.5}
(paper's value).

\begin{center}
\begin{tabular}{lccc}
\toprule
Pangenome metric & Paper & This re-run & Verdict \\
\midrule
Core clusters (all strains, inter-species) & 792--906 & \textbf{909} & \checkmark (within 0.3\% of upper) \\
Pan-genome open at genome \#6            & yes (+553 avg) & yes (+438--698) & \checkmark \\
Strain-specific fraction of pan          & $\sim$65\% & 52.5\% (3{,}123/5{,}946) & $\approx$ (paper used 16 gen.) \\
Core/pan curves monotonic shape          & yes & yes & \checkmark \\
\bottomrule
\end{tabular}
\end{center}

Pan-genome accumulation (average of 30 random orderings):
\begin{center}
\begin{tabular}{lcccccc}
\toprule
genomes added & 1 & 2 & 3 & 4 & 5 & 6 \\
\midrule
pan clusters  & 2459 & 3482 & 4274 & 4926 & 5508 & 5946 \\
core clusters & 2402 & 1327 & 1136 & 1031 & 964  & 909  \\
\bottomrule
\end{tabular}
\end{center}

\section{Layer 2 --- Pathway-content Audit}
Independent scan of \texttt{/function/}, \texttt{/product/}, \texttt{/note/}, and
\texttt{/EC\_number/} fields for the enzymes the paper explicitly calls out:

\begin{center}
\begin{tabularx}{\linewidth}{lXcc}
\toprule
ID & Paper claim & Expected & Observed \\
\midrule
M1 & Methylmalonyl-CoA mutase (5.4.99.16) is only core function & all 6 & 2/2/2/2/2/2 \checkmark \\
M2 & Transaldolase (2.2.1.2) absent only in \emph{P.\ avidum}   & 5/6, PAVI 0 & \textbf{PAVI=0} others $\ge 1$ \checkmark \\
M3 & L-lactate DH (1.1.1.27) in all strains                     & all 6 & 5--6 each \checkmark \\
M4 & Xylose isomerase (5.3.1.5) only in \emph{P.\ acidipropionici} & PAC only & PAC\_4875=2, others=0 \checkmark \\
M5 & Sucrose-6-P hydrolase only in PAC + PPRO                  & 3/3 present, 3/3 absent & exact \checkmark \\
M6 & PFOR/nifJ knockout target in \emph{P.\ freudenreichii}    & PSHE present & CDS.199 / PFREUD\_RS00925 \checkmark \\
\bottomrule
\end{tabularx}
\end{center}

\section{Layer 3 --- Genome-scale FBA Reproduction}
\begin{center}
\begin{tabular}{lcccccc}
\toprule
Model & $\mu$ (default) & $\mu$ (no glc) & Propionate & Open intakes & Vitamins \\
\midrule
PSHE       & 0.786 & \textbf{0.00} & 6.87 & 22 & biotin + pantothenate \\
PAC\_4875  & 0.856 & \textbf{0.00} & 6.49 & 22 & biotin + pantothenate \\
PAC\_55737 & 0.856 & \textbf{0.00} & 6.49 & 22 & biotin + pantothenate \\
PAVI       & 0.854 & 0.00 & 6.52 & 23 & + thiamin \\
PACN       & 0.948 & 0.03 & 9.93 & 29 & + thiamin \\
PPRO       & 1.021 & 0.11 & 4.36 & 28 & + thiamin + riboflavin \\
\bottomrule
\end{tabular}
\end{center}

All six GEMs solve to positive growth under defined media; propionate is the
major fermentation product across all six; and the auxotrophy nesting
(dairy $<$ commensal $<$ opportunist) reproduces exactly.

\section{Claim-by-claim Summary}
\begin{center}
\begin{tabular}{llc}
\toprule
Claim & Layer & Verdict \\
\midrule
C1: All 6 GEMs grow under defined media & FBA & \checkmark \\
C2: Dairy strains glucose-dependent & FBA & \checkmark \\
C3: Propionate the major fermentation product & FBA & \checkmark \\
C4: Auxotrophy hierarchy dairy $<$ commensal $<$ opport. & FBA & \checkmark \\
C5: Vitamin nesting biot/pant $\rightarrow$ +thi $\rightarrow$ +ribo & FBA & \checkmark \\
P1: Core genome 792--906 clusters & Pan-genome & \checkmark (909) \\
P2: Pan-genome open, $\sim$+500/genome & Pan-genome & \checkmark (+438--698) \\
P3: Strain-specific clusters dominate cloud & Pan-genome & \checkmark \\
M1--M6: Diagnostic enzyme presence/absence & Pathway & \checkmark (6/6 exact) \\
\bottomrule
\end{tabular}
\end{center}
\textbf{14/14} numerical / presence-absence claims reproduce on disk.

\section{GENUINE CRITIQUE}
This section documents where the replication is genuinely partial and where the
paper's methodology is thin, in the spirit of a hostile-reader review rather than
a promotional summary.

\subsection{What did \emph{not} independently replicate}
\begin{enumerate}
\item \textbf{Upstream gene-calling step is NOT reproduced from scratch.} The
      paper's protocol begins with a KBase / RAST re-annotation using GLIMMER on
      the 16 raw NCBI GenBank assemblies to standardize gene calls across
      genomes. We did \emph{not} rerun this step; we started from the authors'
      \emph{re-annotated} GenBank files distributed in Supplementary File 4
      (\texttt{Genbank\_files.zip}). Everything downstream --- pan-genome
      clustering, pathway audit, FBA --- is fully independent, but the gene calls
      that seed those analyses are the authors'. This is why coverage is
      \textbf{9/10}, not 10/10, and why the verdict here is \textbf{PARTIAL}
      rather than a full end-to-end replication.
\item \textbf{The KBase/RAST/GLIMMER configuration is underspecified in the
      paper.} Even if we wanted to rerun the annotation, the paper does not pin
      down model versions, exact parameter settings, or the KBase workspace
      snapshot used in 2014--2015. Reproducing this step deterministically
      today would require the frozen KBase pipeline of that era, which is not
      guaranteed to still exist.
\item \textbf{Strain-specific cluster fraction diverges from paper's 65\%.} We
      recover 52.5\% strain-specific clusters across the 6 inter-species reps
      vs.\ the paper's 65\%. The paper's 65\% is computed across all 16 closed
      genomes (which includes 11 \emph{P.\ acnes} strains inflating intra-species
      singletons); ours is 6 inter-species reps. Directionally the same
      conclusion, but the numeric quantity does not match one-for-one.
\end{enumerate}

\subsection{Methodology criticisms of the original paper}
\begin{enumerate}
\item \textbf{OrthoMCL choices are load-bearing but only partially justified.}
      MCL inflation of 1.5, 75\% coverage floor, and $10^{-5}$ e-value are
      standard defaults but the paper does not perform sensitivity analysis
      across neighboring inflation values (e.g.\ 1.2, 2.0, 3.0). The 792--906
      core range the paper reports is itself a symptom of that undocumented
      sensitivity.
\item \textbf{Taxonomy is already outdated.} \emph{Propionibacterium acnes} was
      reassigned to \emph{Cutibacterium acnes} in 2016 (Scholz \& Kilian,
      \emph{Int.\ J.\ Syst.\ Evol.\ Microbiol.}), with \emph{P.\ avidum} and
      \emph{P.\ granulosum} also moving to \emph{Cutibacterium}. A 2020 paper
      still using ``\emph{Propionibacterium acnes}'' throughout its results
      obscures the fact that the ``genus-wide'' pan-genome mixes two now-
      recognized genera.
\item \textbf{Model-quality metrics are not adversarially probed.} The six GEMs
      are declared ``functional'' because they solve for positive growth; but
      the paper does not report standard model-validation metrics (MEMOTE
      score, dead-end reactions, gap-filled reaction fraction, alternate optima
      diversity, essentiality prediction against knockout data).
\item \textbf{Nutritional environment assumptions.} The FBA-derived auxotrophy
      hierarchy depends entirely on the exchange-reaction default bounds coded
      into the GEMs; the paper does not systematically vary media composition.
\item \textbf{No experimental cross-validation of the pathway ``exclusivity''
      claims.} The M4/M5 claims (xylose only in PAC; sucrose-6-P hydrolase only
      in PAC+PPRO) are annotation-based only, not corroborated with growth
      substrate assays in this paper.
\end{enumerate}

\subsection{Where the replication is stronger than the paper}
\begin{itemize}
\item Independent blastp+MCL clustering landing at 909 core (paper 792--906)
      is a near-exact non-bit-identical reproduction, not a re-run of the same
      code.
\item Pathway audit is a fully independent enzyme-level scan of every CDS
      annotation field, not a re-tabulation of the paper's supplementary table.
      All six diagnostic presence/absence patterns match exactly.
\end{itemize}

\section{Verdict}
\textbf{PARTIAL replication (Coverage 9/10, Agreement 10/10 on what was
covered).} Every downstream numerical and presence/absence claim the paper
rests on reproduces on disk. The genuine partial element is that the
upstream KBase/RAST/GLIMMER re-annotation was not re-derived; we used the
authors' deposited standardized GenBank files as the starting point. In
practice this is the maximum reproducibility one can get without a frozen
2014--2015 KBase workspace, and the authors' deposition of the standardized
GenBank files means the paper is \emph{not} reproducibility-blocked --- but
strict end-to-end replication ``from raw NCBI assemblies forward'' remains
untested here.

\section{Artifacts}
\begin{itemize}
\item \texttt{data/genbank/Genbank\_files/*.gbk} --- 6 inter-species rep genomes
\item \texttt{data/proteins/*.faa}, \texttt{all\_proteins.faa}
\item \texttt{data/blast/all\_vs\_all.tsv} --- 414{,}365 hits
\item \texttt{data/pangenome/clusters.txt} --- MCL clusters @ I=1.5
\item \texttt{report/evidence/pangenome.json}, \texttt{pathway\_audit.json},
      \texttt{fba\_replication.json}
\item \texttt{scripts/build\_pangenome.py}, \texttt{extract\_proteins.py},
      \texttt{pathway\_audit.py}, \texttt{fba\_reproduce.py},
      \texttt{inspect\_models.py}
\end{itemize}

\end{document}
